\ifdefined\INCLUDEDFROMTKF\else
\documentclass{article}
\usepackage{amsmath,amssymb,amsthm}
\usepackage{enumitem}
\usepackage{array}
\usepackage{booktabs}

\newtheorem{proposition}{Proposition}
\newtheorem{remark}{Remark}

\newcommand\insrate{\lambda}
\newcommand\delrate{\mu}
\newcommand\ext{r}
\newcommand\notext{\rho}
\newcommand\main{0}
\newcommand\dom{n}
\newcommand\ndom{\mathcal{N}}
\newcommand\nfrag{\mathcal{F}}
\newcommand\domdist{v}
\newcommand\fragdist{w}
\newcommand\classdist{u}
\newcommand\anctok{a}
\newcommand\destok{b}
\newcommand\alphabet{A}
\newcommand\eqm{\pi}
\providecommand\revsub{R}
\providecommand\exch{Q}
\newcommand\evoltime{t}
\newcommand\tkftrans{\tau}
\newcommand\exptrans{\upsilon}
\newcommand\transnest{\chi}
\newcommand\emptyseg{z}
\newcommand\nonemptytrans{T}
\newcommand\slicetrans{U}
\newcommand\betagamma{\beta^\ast}
\newcommand\sta{{\tt S}}
\newcommand\mat{{\tt M}}
\newcommand\ins{{\tt I}}
\newcommand\del{{\tt D}}
\newcommand\fin{{\tt E}}
\newcommand\mnull{{\tt A}}
\newcommand\inull{{\tt B}}
\newcommand\dnull{{\tt C}}
\newcommand\waitm{{\tt V}}
\newcommand\waitd{{\tt W}}
\newcommand\emit{\Omega}
\newcommand\silent{\mathcal{Z}}
\newcommand\effT{\hat{T}}
\newcommand\emprob{e}
\newcommand\pathE{\mathcal{E}}
\newcommand\nullcl{N^\ast}
\newcommand\xstate{{\tt X}}
\newcommand\ystate{{\tt Y}}
\newcommand\ustate{{\tt U}}
\newcommand\vstate{{\tt V}}
\newcommand\srcdom{l}
\newcommand\destdom{m}
\newcommand\frag{f}
\newcommand\srcfrag{f}
\newcommand\destfrag{g}
\newcommand\class{c}
\newcommand\nclasses{\mathcal{C}}

\title{Algebraic Distillation of MixDom}
\author{}
\begin{document}
\maketitle
\fi

\section{Algebraic Distillation of MixDom}
\label{sec:algebraic-distillation}

We investigate whether the distillation of the MixDom model
to order-1 machines (Section~\ref{sec:distill-hmm} and~\ref{sec:order1trans})
can be performed in closed algebraic form, and how the computation scales
with the number of domain types $\ndom$, the number of fragment types
per domain $\nfrag$, and the number of site classes $\nclasses$.
The same algebraic decomposition powers the cherry-count log-likelihood
that the Maraschino fitter (Section~\ref{sec:maraschino-fit}) maximises
under MixDom.

\subsection{Setup}

The MixDom model has the following parameters
(matching Section~\ref{sec:nestpairhmm}):
\begin{itemize}
\item Top-level TKF91: $\insrate_\main, \delrate_\main$ (domain birth/death rates).
\item Domain weights: $\domdist_\dom$, $\sum_\dom \domdist_\dom = 1$.
\item Per-domain TKF91 rates governing the per-domain TKF92 fragment process:
  $\insrate_\dom, \delrate_\dom$ for $\dom = 1, \ldots, \ndom$.
\item Per-domain fragment-type entry distribution:
  $\fragdist_{\dom\frag}$ with $\sum_\frag \fragdist_{\dom\frag} = 1$.
\item Per-domain intra-fragment Markov ext matrix
  $\ext^{(\dom)}_{\srcfrag\destfrag}$, $\nfrag \times \nfrag$, with
  row sums $\leq 1$. The fragment-termination probability is
  $\notext^{(\dom)}_\srcfrag = 1 - \sum_\destfrag \ext^{(\dom)}_{\srcfrag\destfrag}$.
\item Per-(domain, fragment-type) site-class Dirichlet:
  $\classdist_{\dom\frag\class}$ with
  $\sum_\class \classdist_{\dom\frag\class} = 1$.
\item Per-class reversible substitution model:
  $(\exch^{(\class)}, \eqm^{(\class)})$ for $\class = 1, \ldots, \nclasses$;
  rate matrix $\revsub^{(\class)} = \exch^{(\class)} \diag(\eqm^{(\class)})$.
\end{itemize}
The total scalar parameter count is $2 + (\ndom-1) + 2\ndom + \ndom\nfrag^2
+ \ndom(\nfrag-1) + \ndom\nfrag(\nclasses-1)$ plus $\nclasses$ rate matrices.

Write $\kappa_\dom = \insrate_\dom / \delrate_\dom$,
$\alpha_\dom = \alpha(\insrate_\dom, \delrate_\dom, \evoltime)$,
$\beta_\dom = \beta(\insrate_\dom, \delrate_\dom, \evoltime)$,
$\gamma_\dom = \gamma(\insrate_\dom, \delrate_\dom, \evoltime)$,
and similarly $\alpha_\main, \beta_\main, \gamma_\main, \kappa_\main$
for the top-level parameters.

\subsection{Class-mixture emissions}
\label{sec:class-mixture-emissions}

In MixDom the per-(domain, fragment-type) site-class mixture
appears in every emission probability of the collapsed Pair HMM.
Define the per-class match emission as
\[
P^{(\class)}(\anctok,\destok\mid\evoltime)
= \eqm^{(\class)}_\anctok\, \exp(\revsub^{(\class)}\evoltime)_{\anctok\destok}.
\]
The model's per-(domain, fragment-type) emission factors at evolutionary time $\evoltime$ are:
\begin{align*}
\phi^{\mat}_{\dom\frag}(\anctok, \destok)
  &\equiv \emprob_{\mat\mat_{\dom\frag}}(\anctok,\destok)
  = \sum_{\class=1}^{\nclasses} \classdist_{\dom\frag\class}\,
    P^{(\class)}(\anctok,\destok\mid\evoltime), \\
\phi^{\ins}_{\dom\frag}(\destok)
  &\equiv \emprob_{\mat\ins_{\dom\frag}}(\destok)
       = \emprob_{\ins\ins_{\dom\frag}}(\destok)
  = \sum_\class \classdist_{\dom\frag\class}\, \eqm^{(\class)}_\destok, \\
\phi^{\del}_{\dom\frag}(\anctok)
  &\equiv \emprob_{\mat\del_{\dom\frag}}(\anctok)
       = \emprob_{\del\del_{\dom\frag}}(\anctok)
  = \sum_\class \classdist_{\dom\frag\class}\, \eqm^{(\class)}_\anctok.
\end{align*}
Each of these is a $\nclasses$-term mixture over the per-class GTRs.
The match emission $\phi^{\mat}_{\dom\frag}$ inherits the time
dependence from the per-class transition kernels $P^{(\class)}(\evoltime)$;
the singlet emissions $\phi^{\ins}_{\dom\frag}$, $\phi^{\del}_{\dom\frag}$
are time-independent.

\begin{remark}[Bilinear emission tensor]
The match emission factors as
\[
\phi^{\mat}_{\dom\frag}(\anctok,\destok)
= \sum_{\class} \classdist_{\dom\frag\class}\, \eqm^{(\class)}_\anctok\,
  \exp(\revsub^{(\class)}\evoltime)_{\anctok\destok}
\]
which is rank-$\nclasses$ in the latent class index, and the resulting
$|\alphabet| \times |\alphabet|$ emission tensor at $(\dom,\frag)$ has
rank at most $\nclasses$. The single-class case $\classdist = \delta_{\dom\class}$
makes this rank-1 per domain (a single GTR per domain). The general parameterisation strictly
generalises by sharing the GTR pool across (domain, fragment-type)
contexts and allowing any soft mixture.
\end{remark}

\subsection{Single HMM Distillation}

\subsubsection{State space}

The MixDom Singlet HMM generates sequences from the stationary
distribution. Its collapsed state space (Section~\ref{sec:nestpairhmm})
is $\{ \sta, \fin \} \cup \{ \ins_{\dom\frag} : \dom \in \ndom,
\frag \in \nfrag \}$, so there are $\ndom\nfrag$ emitting states with
emissions $\phi^\ins_{\dom\frag}(\anctok)$.

Within a domain, fragment continuation is governed by the intra-fragment
Markov ext matrix $\ext^{(\dom)}_{\srcfrag\destfrag}$ and the
new-fragment $\kappa_\dom \fragdist_{\dom,\destfrag}$ branch.
Adapting the formulae in Section~\ref{sec:nestpairhmm} for the singlet,
the effective transition matrix between emitting states
$\ins_{\srcdom\srcfrag} \to \ins_{\destdom\destfrag}$ has entries
\[
\effT_{\ins_{\srcdom\srcfrag},\ins_{\destdom\destfrag}}
= \ext^{(\srcdom)}_{\srcfrag\destfrag}\delta_{\srcdom\destdom}
+ \notext^{(\srcdom)}_\srcfrag \kappa_\srcdom \fragdist_{\srcdom\destfrag}\delta_{\srcdom\destdom}
+ \frac{\notext^{(\srcdom)}_\srcfrag (1-\kappa_\srcdom)\,
       \kappa_\main \domdist_\destdom \kappa_\destdom \fragdist_{\destdom\destfrag}}
       {1 - \kappa_\main \emptyseg_0}
\]
where $\emptyseg_0 = \sum_\dom \domdist_\dom (1-\kappa_\dom)$.
The path-sum matrix $(I - \effT_{\emit,\emit})^{-1}$ is an
$\ndom\nfrag \times \ndom\nfrag$ inversion in general.

\subsubsection{Adjacency frequencies}

The Singlet adjacency frequency is
\begin{equation}\label{eq:single-adj}
f(\anctok, \destok)
= \sum_{(\srcdom,\srcfrag),(\destdom,\destfrag)}
  W_{(\srcdom,\srcfrag),(\destdom,\destfrag)}\,
  \phi^\ins_{\srcdom\srcfrag}(\anctok)\,
  \phi^\ins_{\destdom\destfrag}(\destok)
\end{equation}
where the structural weights
$W_{(\srcdom,\srcfrag),(\destdom,\destfrag)}
= L_{(\srcdom,\srcfrag)}\, \effT_{(\srcdom,\srcfrag)(\destdom,\destfrag)}\,
R_{(\destdom,\destfrag)}$ are character-independent, with
$L_{(\dom,\frag)} = \sum_i \pi_i \effT^\ast_{i,(\dom,\frag)}$ and
$R_{(\dom,\frag)} = \sum_j \effT^\ast_{(\dom,\frag),j}$.
Substituting the class-mixture emission yields a bilinear sum over
both the (domain, fragment-type) latent and the site-class latent:
\[
f(\anctok, \destok)
= \sum_{\srcdom,\srcfrag,\destdom,\destfrag,\class_1,\class_2}
  W_{(\srcdom,\srcfrag),(\destdom,\destfrag)}
  \classdist_{\srcdom\srcfrag\class_1}
  \classdist_{\destdom\destfrag\class_2}\,
  \eqm^{(\class_1)}_\anctok\, \eqm^{(\class_2)}_\destok.
\]

The order-1 HMM transition is therefore
\[
P(\destok \mid \anctok)
= \frac{f(\anctok, \destok)}{\sum_{\destok'} f(\anctok, \destok')}
= \frac{
    \sum_{\class_2}
    \big(\sum_{\srcdom,\srcfrag,\destdom,\destfrag,\class_1}
         W \classdist \classdist\, \eqm^{(\class_1)}_\anctok\big)
    \eqm^{(\class_2)}_\destok
  }{
    \sum_{\srcdom,\srcfrag,\class_1} \eqm^{(\class_1)}_\anctok
    \sum_{\destdom,\destfrag,\class_2}
      W_{(\srcdom,\srcfrag),(\destdom,\destfrag)}
      \classdist_{\srcdom\srcfrag\class_1}
      \classdist_{\destdom\destfrag\class_2}
  }.
\]

\begin{remark}[Non-trivial character correlations]
When the per-class equilibria $\eqm^{(\class)}$ differ, the previous
character $\anctok$ carries information about the latent class $\class_1$
of the source state via the Bayesian posterior
$P(\class_1 \mid \anctok) \propto \eqm^{(\class_1)}_\anctok\,
\sum_{\srcdom\srcfrag,\destdom\destfrag,\class_2}
W \classdist \classdist$,
and consequently about which (domain, fragment-type) is likely to be
generating this region. The order-1 HMM therefore has genuinely
$\anctok$-dependent transitions: a mixture of $\eqm^{(\class_2)}$ tilted
by the joint posterior over latent state and source class.
\end{remark}

\begin{remark}[Closed form]
All quantities in~\eqref{eq:single-adj} are rational functions of the
model parameters: the $\ndom\nfrag \times \ndom\nfrag$ matrix
$(I - \effT)^{-1}$ has entries that are ratios of polynomials in
$(\ext^{(\dom)}_{\srcfrag\destfrag}, \kappa_\dom, \kappa_\main, \beta_\main,
\fragdist_{\dom\frag}, \domdist_\dom)$. The class mixture enters
linearly through $\classdist_{\dom\frag\class}$, and the per-class GTRs
appear only inside $\eqm^{(\class)}, P^{(\class)}(\evoltime)$.
\end{remark}


\subsection{Pair HMM Distillation}

\subsubsection{Emitting state space}

The collapsed MixDom Pair HMM with $\ndom$ domain types and $\nfrag$
fragment types per domain has $5\ndom\nfrag + 2$ states:
$\{ \sta\sta, \fin\fin \}$ and
$\{ \mat\mat_{\dom\frag}, \mat\ins_{\dom\frag}, \mat\del_{\dom\frag},
\ins\ins_{\dom\frag}, \del\del_{\dom\frag}
: \dom \in \ndom, \frag \in \nfrag \}$.

Group the $5\ndom\nfrag$ emitting states by emission type:
match $\mathcal{M} = \{\mat\mat_{\dom\frag}\}$,
insert $\mathcal{I} = \{\mat\ins_{\dom\frag}, \ins\ins_{\dom\frag}\}$,
delete $\mathcal{D} = \{\mat\del_{\dom\frag}, \del\del_{\dom\frag}\}$.

\subsubsection{Class-mixture emissions in the Pair HMM}

The per-state emissions at $(d,f)$ are the class-mixture emissions
$\phi^{\mat,\ins,\del}_{\dom\frag}$ defined in
Section~\ref{sec:class-mixture-emissions}. Two states with the same
$(d,f)$ but different top-level types share the same equilibrium
mixture, but the match state combines this with the time-evolved
class-mixture transition kernel.

The per-(domain, fragment-type) emission factors do \emph{not} lift
cleanly out of the structural sum; however, they have low rank in the
class index $\class$ (Section~\ref{sec:class-mixture-emissions}).
Within a fixed $(d,f)$, all match states emit from
$\phi^{\mat}_{\dom\frag}$ and all insert/delete singlets from
$\phi^{\ins}_{\dom\frag}$ / $\phi^{\del}_{\dom\frag}$.

\subsubsection{Per-(domain, fragment)-pair structural weights}

The pair adjacency frequencies decompose as sums over
(domain, fragment-type) pairs.
For match-to-match, writing
$X = \anctok, Y = \destok, X' = \anctok', Y' = \destok'$:
\begin{equation}\label{eq:pair-adj}
f^{\mathcal{M}\mathcal{M}}(X,Y,X',Y')
= \sum_{(\dom_1,\frag_1),(\dom_2,\frag_2)}
  C^{\mathcal{M}\mathcal{M}}_{(\dom_1,\frag_1),(\dom_2,\frag_2)}\,
  \phi^{\mat}_{\dom_1\frag_1}(X,Y)\,
  \phi^{\mat}_{\dom_2\frag_2}(X',Y')
\end{equation}
where
\[
C^{\mathcal{M}\mathcal{M}}_{(\dom_1,\frag_1),(\dom_2,\frag_2)}
= \sum_{s_1 \in \mathcal{M}_{\dom_1\frag_1}} \sum_{s_2 \in \mathcal{M}_{\dom_2\frag_2}}
  L_{s_1}\, \effT_{s_1 s_2}\, R_{s_2}
\]
are structural weights indexed by (domain, fragment-type) pairs (not
just emission types). Similarly for other adjacency types
($\mathcal{M}\mathcal{I}$, $\mathcal{M}\mathcal{D}$, etc.), with
the appropriate per-(domain, fragment-type) emissions.

Each adjacency frequency is therefore a sum of $(\ndom\nfrag)^2$
structural-weight terms, each multiplied by per-(domain, fragment-type)
emission factors that are themselves rank-$\nclasses$ class mixtures.

\subsubsection{Context dependence}

In the order-1 WFST, the wait-after-match/insert transition involves
\begin{align*}
&f^{\cdot\mathcal{M}}(X,Y,\anctok',\destok')
= \sum_{(\dom_1,\frag_1),(\dom_2,\frag_2)}
  \Big( C^{\mathcal{M}\mathcal{M}}_{(\dom_1,\frag_1),(\dom_2,\frag_2)}\,
        \phi^{\mat}_{\dom_1\frag_1}(X,Y) \\
&\phantom{= \sum_{(\dom_1,\frag_1),(\dom_2,\frag_2)} \Big(}
      + C^{\mathcal{I}\mathcal{M}}_{(\dom_1,\frag_1),(\dom_2,\frag_2)}\,
        \phi^{\ins}_{\dom_1\frag_1}(Y) \Big)
  \phi^{\mat}_{\dom_2\frag_2}(\anctok',\destok')
\end{align*}

After normalization, $p_{\waitm\mat}(X,Y,\anctok',\destok')$ depends on
$(X,Y)$ through the joint posterior over (domain, fragment-type, class)
of the previous emission. Specifically, the context enters through the
$\ndom\nfrag$-dimensional vector
\[
\boldsymbol{\rho}(X,Y) = \Big(
  \sum_{(\dom_2,\frag_2)} C^{\mathcal{M}\mathcal{M}}_{(1,1),(\dom_2,\frag_2)}\,
  \phi^{\mat}_{1,1}(X,Y)
  + \sum_{(\dom_2,\frag_2)} C^{\mathcal{I}\mathcal{M}}_{(1,1),(\dom_2,\frag_2)}\,
  \phi^{\ins}_{1,1}(Y), \ldots \Big)
\]
which captures the relative likelihood of each source (domain, fragment-type).

\begin{remark}[Richer context than per-domain shared-$\eqm$ case]
With class-mixture emissions, the WFST transition probabilities depend
on context $(X,Y)$ through a joint (domain, fragment-type, class)
posterior that is sensitive to \emph{both} $X$ and $Y$ independently---not
just through a single scalar ratio. The right-side character dependence
($\anctok', \destok'$) also varies by (domain, fragment-type, class):
the next emission is a mixture of
$\eqm^{(\class_2)} P^{(\class_2)}(\evoltime)$ weighted by the structural
weights, the source posterior, and the destination $\classdist$.
Despite this richer structure, all quantities are closed-form rational
functions of the model parameters (composed with $\nclasses$ rate-matrix
exponentials).
\end{remark}


\subsection{Block Structure and Matrix Inversions}

\subsubsection{\texorpdfstring{Top-level null closure (always $3 \times 3$)}{Top-level null closure (always 3 x 3)}}

The null states $\mnull, \inull, \dnull$ yield a $3 \times 3$ submatrix
with entries depending on
\begin{align*}
\emptyseg_0 &= \sum_\dom \domdist_\dom(1 - \kappa_\dom) \\
\emptyseg_\evoltime &= \sum_\dom \domdist_\dom(1 - \kappa_\dom)(1 - \beta_\dom)
\end{align*}
The null closure is a $3 \times 3$ inversion with closed-form
determinant, independent of $\ndom, \nfrag, \nclasses$ (the dependence
on $\ndom$ enters only through these two scalar sums; $\nfrag$ and
$\nclasses$ do not appear at all because empty domains are independent
of fragment and class structure inside the domain).

\subsubsection{Block-diagonal-plus-low-rank decomposition}

The $5\ndom\nfrag \times 5\ndom\nfrag$ effective transition matrix
between emitting states decomposes as
\[
\effT_{\emit,\emit}
= \underbrace{\mathrm{diag}(D_1, \ldots, D_\ndom)}_{\text{intra-domain}}
+ \underbrace{E\,\nonemptytrans_\bullet\, S^\top}_{\text{inter-domain (rank } \leq 3\text{)}}
\]
where each $D_\dom$ is a $5\nfrag \times 5\nfrag$ within-domain block
that combines the intra-fragment Markov ext matrix
$\ext^{(\dom)}_{\srcfrag\destfrag}$ (acting on the fragment-type axis)
with the same-domain new-fragment TKF92 transitions
($\notext^{(\dom)}_\srcfrag \tkftrans^{(\dom)}_{\xstate\ystate}
\fragdist_{\dom,\destfrag}$ for M-type entries,
$\notext^{(\dom)}_\srcfrag \kappa_\dom \fragdist_{\dom,\destfrag}$ for
I/D-type entries),
$E \in \mathbb{R}^{5\ndom\nfrag \times 3}$ stacks per-(domain,fragment-type)
exit vectors projected to top-level types $\{\mat,\ins,\del\}$ via
$(\notext^{(\dom)}_\srcfrag\, \notkappa_\dom)$ for the I/D singlet rows
and $(\notext^{(\dom)}_\srcfrag\, \tkftrans^{(\dom)}_{\xstate\fin})$ for
the matched-domain rows,
$S \in \mathbb{R}^{5\ndom\nfrag \times 3}$ stacks per-(domain,fragment-type)
start vectors with $(1-\emptyseg_\evoltime)^{-1}\domdist_\destdom\,
\tkftrans^{(\destdom)}_{\sta\ystate}\fragdist_{\destdom\destfrag}$ for
M-type entries and $(1-\emptyseg_0)^{-1}\domdist_\destdom\, \kappa_\destdom\,
\fragdist_{\destdom\destfrag}$ for I/D-type entries (these
$1/(1-\emptyseg)$ factors are the domenter normalisation that conditions
on the destination domain being non-empty),
and $\nonemptytrans_\bullet$ is the relevant $5 \times 5$ submatrix of
the null-eliminated top-level matrix.

Note that $D_{\dom_1} \neq D_{\dom_2}$ in general (different
$\ext^{(\dom)}_{\srcfrag\destfrag}, \fragdist_{\dom\frag}, \insrate_\dom,
\delrate_\dom$). For $\nfrag > 1$, $D_\dom$ has no further
emission-type block-diagonal sub-structure within the matched-domain
block (the intra-fragment Markov ext couples
$\{\mat\mat, \mat\ins, \mat\del\}$ across fragment-types); however, it
admits a Kronecker-plus-low-rank decomposition that preserves the
linear-in-$\ndom\nfrag$ scaling
(Section~\ref{sec:within-domain-mixdom}). The diagonal-ext baseline
(\texttt{--freeze-offdiag-ext} in the fitter) recovers the diagonal-extension special case
block structure as a special case.

\begin{remark}[Domenter normalization]
The $\transnest$ matrix (Section~\ref{sec:nestpairhmm},
Equation~\ref{eq:mixdom_transitions}) includes factors
$(1-\emptyseg_\evoltime)^{-1}$ for M-type destinations and
$(1-\emptyseg_0)^{-1}$ for I/D-type destinations, conditioning domain
entry on the domain being non-empty. These factors cancel the
corresponding $(1-\emptyseg_\evoltime)$ and $(1-\emptyseg_0)$ column
scaling in $\nonemptytrans_\bullet$ (which arises from the $\exptrans$
null-elimination Schur complement). In the distillation, the start
vectors $\mathbf{s}_{\dom,\frag}$ carry the domenter factors while
$\nonemptytrans_\bullet$ retains the column factors; the product
$\nonemptytrans_\bullet \cdot \mathbf{s}_{\dom,\frag}$ thus reproduces
the correct $\transnest$ entries.
\end{remark}

\subsubsection{Woodbury identity}

Writing $G_\dom = (I-D_\dom)^{-1}$ (each $5\nfrag \times 5\nfrag$),
$E = (\mathbf{e}_{\dom,\frag})$ (exit vectors, $5\ndom\nfrag \times 3$),
$S = (\mathbf{s}_{\dom,\frag})$ (start vectors, $5\ndom\nfrag \times 3$),
and $\nonemptytrans_\text{mid}$ for the $3\times 3$ subblock of
$\nonemptytrans_\bullet$ restricted to $\{\mat,\ins,\del\}$ (the only
top-level types that couple across domains), the path-sum matrix is
\[
(I - \effT)^{-1}
= G + G\,E\,\nonemptytrans_\text{mid}
  \bigl(I_3 - S^\top G\,E\,\nonemptytrans_\text{mid}\bigr)^{-1}
  S^\top G
\]
using the push-through form of the Woodbury identity, where
$G = \mathrm{diag}(G_1, \ldots, G_\ndom)$.
This avoids inverting $\nonemptytrans_\text{mid}$, which is singular
(the TKF91 $\mat$ and $\ins$ rows are identical, so
$\nonemptytrans_\text{mid}$ has rank~2).

\begin{proposition}[Closed-form distillation for MixDom]
All adjacency frequencies, and hence all order-1 HMM and WFST
parameters, are closed-form rational functions of the model parameters
(composed with the matrix exponentials $\exp(\revsub^{(\class)}\evoltime)$),
for any finite $\ndom, \nfrag, \nclasses$.
The computation requires:
\begin{enumerate}
\item A $3 \times 3$ inversion for the top-level null closure.
\item $\ndom$ within-domain inversions $(I - D_\dom)^{-1}$, each of size
  $5\nfrag \times 5\nfrag$. By the Kronecker-plus-rank-3 decomposition
  of $D_\dom$ (Section~\ref{sec:within-domain-mixdom}), this reduces
  to one $\nfrag \times \nfrag$ inversion of
  $(I - \ext^{(\dom)})$ plus a $3 \times 3$ inner Woodbury kernel per
  domain. For $\nfrag = 1$ the whole within-domain block collapses to a
  $3 \times 3$ adjugate plus two scalar inversions
  (Section~\ref{sec:within-domain-inversion}).
\item One $3 \times 3$ outer Woodbury correction for inter-domain coupling.
\item Summation of $(\ndom\nfrag)^2$ (domain, fragment-type)-pair terms
  per adjacency entry, with each term involving a class-mixture
  emission factor of length $\nclasses$.
\end{enumerate}
Every step is a rational function of the model parameters: no
numerical iteration is required to evaluate the closed form.
\end{proposition}

\begin{remark}[Contrast with shared-emission case]
If all $(d,f)$ shared a single class mixture (e.g.\
$\classdist_{\dom\frag\class} = u_\class$ uniform), the adjacency
frequencies would still factor as in~\eqref{eq:pair-adj}, but the
emission tensor would have rank-1 structure across $(d,f)$, collapsing
the WFST context dependence. With per-(domain, fragment-type) class
mixtures, the structural constants are full
$(\ndom\nfrag) \times (\ndom\nfrag)$ matrices indexed by
(domain, fragment-type) pairs, the emission tensors are
class-rank-$\nclasses$, and the context dependence is genuinely
multi-dimensional.
\end{remark}


\subsection{Within-Domain Inversion: closed form}
\label{sec:within-domain-mixdom}

For general $\nfrag$, the within-domain block $D_\dom$ has a
Kronecker-plus-rank-3 structure that yields a closed-form inverse
without ever inverting a $5\nfrag \times 5\nfrag$ matrix directly.

\subsubsection{Kronecker-plus-rank-3 decomposition of the matched-domain block}

The matched-domain inner block (acting on
$\{\mat\mat, \mat\ins, \mat\del\}$ at each fragment-type, total
$3\nfrag$ states) combines intra-fragment Markov continuation with
fragment termination plus a TKF92 same-domain new-fragment branch:
\begin{equation}\label{eq:mixdom-D-match}
D_\dom^{\text{match}}[(\xstate, \frag), (\ystate, \destfrag)]
= \delta_{\xstate\ystate}\, \ext^{(\dom)}_{\frag\destfrag}
+ \notext^{(\dom)}_\frag \cdot
  \tkftrans^{(\dom)}_{\xstate\ystate} \cdot
  \fragdist_{\dom,\destfrag}.
\end{equation}
The first term is $I_{3} \otimes \ext^{(\dom)}$ (the
intra-fragment Markov chain, identical at every emission type);
the second is
$\tkftrans^{(\dom)}_{3\times 3} \otimes (\notext^{(\dom)}\, \fragdist_\dom^\top)$
where $\notext^{(\dom)}\, \fragdist_\dom^\top$ is the rank-1
$\nfrag \times \nfrag$ outer product of fragment-termination
probabilities and fragment-entry weights.
The whole rank-3 correction is therefore of rank $\leq 3$ in the joint
$(\xstate,\frag)$ space (one factor of 3 from
$\tkftrans^{(\dom)}_{3\times 3}$, and rank 1 in the fragment slot
from the outer product). Writing
\[
D_\dom^{\text{match}} = I_{3} \otimes \ext^{(\dom)}
+ U_\dom\, V_\dom^\top,
\quad
U_\dom = \tkftrans^{(\dom)}_{3\times 3} \otimes \notext^{(\dom)},
\quad
V_\dom = I_{3} \otimes \fragdist_\dom,
\]
$U_\dom$ and $V_\dom$ are $3\nfrag \times 3$ matrices.

\subsubsection{Inverse via inner Woodbury}

Since $I_{3} \otimes \ext^{(\dom)}$ commutes with itself,
$I - I_3 \otimes \ext^{(\dom)} = I_3 \otimes (I_\nfrag - \ext^{(\dom)})$,
which inverts blockwise: $(I_3 \otimes (I_\nfrag - \ext^{(\dom)}))^{-1}
= I_3 \otimes (I_\nfrag - \ext^{(\dom)})^{-1}$.
Note $(I_\nfrag - \ext^{(\dom)})^{-1}$ is a single
$\nfrag \times \nfrag$ inversion per domain.
Applying the Sherman--Morrison--Woodbury identity to the rank-3
correction gives
\begin{equation}\label{eq:mixdom-Dmatch-inverse}
\begin{aligned}
(I - D_\dom^{\text{match}})^{-1}
&= G_0
+ G_0\, U_\dom\, K_\dom^{\text{inner}}\, V_\dom^\top G_0, \\
G_0 &\equiv I_{3} \otimes (I_\nfrag - \ext^{(\dom)})^{-1}, \\
K_\dom^{\text{inner}}
&\equiv (I_3 - V_\dom^\top G_0\, U_\dom)^{-1}.
\end{aligned}
\end{equation}
$K_\dom^{\text{inner}}$ is a $3 \times 3$ matrix whose entries are
rational functions of $\fragdist_\dom$, $\ext^{(\dom)}$,
$\tkftrans^{(\dom)}$ via a single $\nfrag$-fold inner product
$\fragdist_\dom^\top (I - \ext^{(\dom)})^{-1} \notext^{(\dom)}$
times $\tkftrans^{(\dom)}_{3\times 3}$.

The two singlet-domain blocks act on $\nfrag$ states each
($\{\ins\ins_{\dom\frag} : \frag\}$ and
$\{\del\del_{\dom\frag} : \frag\}$):
\[
d_\dom^{\text{ins}}[\frag,\destfrag]
= d_\dom^{\text{del}}[\frag,\destfrag]
= \ext^{(\dom)}_{\frag\destfrag} + \notext^{(\dom)}_\frag\,
  \kappa_\dom\, \fragdist_{\dom,\destfrag}.
\]
Their inverse follows the same pattern: rank-1 correction to
$(I - \ext^{(\dom)})$, giving
\[
(I - d_\dom^{\text{ins}})^{-1}
= (I - \ext^{(\dom)})^{-1}
+ (I - \ext^{(\dom)})^{-1} \notext^{(\dom)}\,
  k_\dom^{\text{ins}}\,
  \fragdist_\dom^\top (I - \ext^{(\dom)})^{-1},
\]
with scalar inner kernel
$k_\dom^{\text{ins}}
= \kappa_\dom \, [1 - \kappa_\dom \fragdist_\dom^\top
(I - \ext^{(\dom)})^{-1} \notext^{(\dom)}]^{-1}$,
and identically for delete.

Every entry of $(I - D_\dom)^{-1}$ is therefore a rational function of
$(\ext^{(\dom)}, \fragdist_\dom, \alpha_\dom, \beta_\dom, \gamma_\dom,
\kappa_\dom)$, computed at cost $O(\nfrag^3)$ per domain (the inversion
of $I_\nfrag - \ext^{(\dom)}$) plus $O(\nfrag^2)$ for the singlet
blocks and $O(1)$ for the inner $3 \times 3$ Woodbury kernel.
No numerical iteration is required; the inversion can be carried out
analytically (for small $\nfrag$, by Cramer's rule on
$I_\nfrag - \ext^{(\dom)}$) or directly evaluated for larger $\nfrag$.

\begin{remark}[Compactness]
For small $\nfrag$ the inverse $(I_\nfrag - \ext^{(\dom)})^{-1}$ has a
compact symbolic form (e.g.\ at $\nfrag = 2$ it is a $2 \times 2$
adjugate over a scalar determinant). The further reductions of the
matched-domain block via~\eqref{eq:mixdom-Dmatch-inverse} and the
singlet blocks above involve only $3 \times 3$ and scalar inner
inversions, so the overall within-domain inverse can be
written down symbolically without ever exceeding a $\nfrag \times
\nfrag$ inversion. The $\nfrag = 1$ scalar-extension special case
collapses $(I_\nfrag - \ext^{(\dom)})^{-1}$ to the scalar
$1/(1 - \ext_\dom)$ and recovers the more compact
$3 \times 3$ adjugate of Section~\ref{sec:within-domain-inversion}
below.
\end{remark}

\subsection{\texorpdfstring{Within-Domain Inversion: $\nfrag = 1$ closed form}{Within-Domain Inversion: F = 1 closed form}}
\label{sec:within-domain-inversion}

When $\nfrag = 1$ (scalar self-extension
$\ext_\dom \equiv \ext^{(\dom)}_{11}$, no off-diagonal Markov coupling),
$D_\dom$ collapses to $5 \times 5$ with the additional block structure
described below; we record this special case because it admits a
particularly compact algebraic form.

\subsubsection{\texorpdfstring{Block decomposition: $3 \times 3 + 1 + 1$}{Block decomposition: 3 x 3 + 1 + 1}}

In the $\nfrag = 1$ limit, the five emitting state types
$\{\mat\mat, \mat\ins, \mat\del, \ins\ins, \del\del\}$
cannot transition between top-level types within a domain:
if the domain is inserted (top-level $\ins$), all emissions stay in
$\ins\ins$ until the domain ends; similarly for deleted domains
($\del\del$). Therefore $D_\dom$ is block-diagonal:
\[
D_\dom = \begin{pmatrix} D_\dom^{\text{match}} & 0 & 0 \\ 0 & d_\dom^{\text{ins}} & 0 \\ 0 & 0 & d_\dom^{\text{del}} \end{pmatrix}
\]
where $D_\dom^{\text{match}}$ is $3 \times 3$ (for matched-domain states
$\mat\mat, \mat\ins, \mat\del$) and the singlet-domain states have
scalar self-loops:
\[
d_\dom^{\text{ins}} = d_\dom^{\text{del}} = \ext_\dom + (1 - \ext_\dom)\kappa_\dom
\]

The scalar inversions are trivial:
$(1 - d_\dom^{\text{ins}})^{-1} = [(1-\ext_\dom)(1-\kappa_\dom)]^{-1}$.

\subsubsection{\texorpdfstring{The $3\times 3$ matched-domain block}{The 3 x 3 matched-domain block}}

The matched-domain block combines fragment extension (self-loop at rate
$\ext_\dom$) with intra-domain TKF transitions:
\[
D_\dom^{\text{match}} = \ext_\dom I_3 + (1 - \ext_\dom)\, \tkftrans_{\mat\ins\del,\mat\ins\del}^{(\dom)}
\]
where $\tkftrans_{\mat\ins\del,\mat\ins\del}^{(\dom)}$ is the
$3 \times 3$ submatrix of $\tkftrans^{(\dom)}$ restricted to rows and
columns $\mat, \ins, \del$.

A key property: in the TKF transition matrix, rows $\mat$ and $\ins$
are identical. Defining shorthand (suppressing domain subscript $\dom$):
\begin{align*}
\mathfrak{a} &= (1-\beta)\kappa\alpha, &
\mathfrak{b} &= \beta, &
\mathfrak{c} &= (1-\beta)\kappa(1-\alpha) \\
\mathfrak{d} &= (1-\gamma)\kappa\alpha, &
\mathfrak{g} &= \gamma, &
\mathfrak{h} &= (1-\gamma)\kappa(1-\alpha)
\end{align*}
the $3 \times 3$ TKF submatrix is
\[
\tkftrans_{\mat\ins\del,\mat\ins\del}^{(\dom)} =
\begin{pmatrix}
\mathfrak{a} & \mathfrak{b} & \mathfrak{c} \\
\mathfrak{a} & \mathfrak{b} & \mathfrak{c} \\
\mathfrak{d} & \mathfrak{g} & \mathfrak{h}
\end{pmatrix}
\]
with row sums $\mathfrak{a}+\mathfrak{b}+\mathfrak{c} = (1-\beta)\kappa + \beta$ for the first two rows
and $\mathfrak{d}+\mathfrak{g}+\mathfrak{h} = (1-\gamma)\kappa + \gamma$ for the third.

\subsubsection{Factoring out the extension rate}

Since $I - D^{\text{match}} = I - \ext I - (1-\ext)\tkftrans_{3\times 3}
= (1-\ext)(I - \tkftrans_{3\times 3})$,
the extension rate factors out as a scalar:
\[
(I - D_\dom^{\text{match}})^{-1} = \frac{1}{1-\ext_\dom}\, P_\dom^{-1}
\]
where
\[
P_\dom \equiv I - \tkftrans_{\mat\ins\del,\mat\ins\del}^{(\dom)} =
\begin{pmatrix}
1-\mathfrak{a} & -\mathfrak{b} & -\mathfrak{c} \\
-\mathfrak{a} & 1-\mathfrak{b} & -\mathfrak{c} \\
-\mathfrak{d} & -\mathfrak{g} & 1-\mathfrak{h}
\end{pmatrix}
\]
with row sums $(1-\beta)(1-\kappa)$, $(1-\beta)(1-\kappa)$, $(1-\gamma)(1-\kappa)$
(the domain-exit probabilities from each inner state).

\subsubsection{\texorpdfstring{Determinant of $P_\dom$}{Determinant of P\_n}}

Since rows 1 and 2 of $\tkftrans_{3\times 3}$ are identical,
subtracting row~2 from row~1 in $P$ yields $(1, -1, 0)$.
Expanding the determinant along this simplified row gives
\begin{align*}
\det(P_\dom) &= (1-\mathfrak{b})(1-\mathfrak{h}) - \mathfrak{c}\mathfrak{g}
   - \mathfrak{a}(1-\mathfrak{h}) - \mathfrak{c}\mathfrak{d} \\
&= (1-\beta_\dom)(1-\kappa_\dom)
\end{align*}
This factors cleanly: since $\mathfrak{a}+\mathfrak{b}+\mathfrak{c}$ and
$\mathfrak{d}+\mathfrak{g}+\mathfrak{h}$ are the row sums of
$\tkftrans_{3\times 3}$, we have
\[
\boxed{\det(P_\dom) = (1-\beta_\dom)(1-\kappa_\dom)}
\]
The determinant is the product of the M/I-row exit probability $(1-\beta)$
and the stationary emptiness probability $(1-\kappa)$.

\subsubsection{\texorpdfstring{Inverse of $P_\dom$}{Inverse of P\_n}}

The adjugate $\text{adj}(P_\dom)$ has entries (again suppressing domain subscripts):
\[
\text{adj}(P) =
\begin{pmatrix}
(1-\beta)(1-\kappa(1-\alpha))
  & \beta + \kappa(1-\alpha)(\gamma-\beta)
  & (1-\beta)\kappa(1-\alpha) \\[4pt]
(1-\beta)\kappa\alpha
  & 1 - (1-\beta)\kappa\alpha - (1-\gamma)\kappa(1-\alpha)
  & (1-\beta)\kappa(1-\alpha) \\[4pt]
(1-\beta)\kappa\alpha
  & \gamma + \kappa\alpha(\beta-\gamma)
  & (1-\beta)(1-\kappa\alpha)
\end{pmatrix}
\]
Note the symmetries: $\text{adj}(P)_{21} = \text{adj}(P)_{31}$
and $\text{adj}(P)_{13} = \text{adj}(P)_{23}$
(reflecting the identical rows of $\tkftrans_{3\times 3}$),
while $\text{adj}(P)_{11} - \text{adj}(P)_{33} = (1-\beta)\kappa(2\alpha-1)$,
which vanishes only when $\alpha = \tfrac{1}{2}$.

The full inverse is
\[
P_\dom^{-1} = \frac{\text{adj}(P_\dom)}{(1-\beta_\dom)(1-\kappa_\dom)}
\]
and therefore
\[
(I - D_\dom^{\text{match}})^{-1} = \frac{\text{adj}(P_\dom)}{(1-\ext_\dom)(1-\beta_\dom)(1-\kappa_\dom)}
\]

\begin{remark}[Compactness in the $\nfrag = 1$ limit]
Each entry of $(I - D_\dom^{\text{match}})^{-1}$ is a ratio of a polynomial
with 1--3 terms (numerator, from the adjugate) over a polynomial with 2--3 terms
(denominator $(1-\ext)(1-\beta)(1-\kappa)$). These are genuinely compact
closed-form expressions in the TKF parameters
$(\alpha_\dom, \beta_\dom, \gamma_\dom, \kappa_\dom, \ext_\dom)$.
For $\nfrag > 1$ the inversion is still closed-form (a rational
function in $\ext^{(\dom)}, \fragdist_\dom, \alpha_\dom, \beta_\dom,
\gamma_\dom, \kappa_\dom$), via the Kronecker-plus-rank-3 Woodbury
decomposition of Section~\ref{sec:within-domain-mixdom}: the only
inversion of dimension exceeding $3 \times 3$ is a single
$\nfrag \times \nfrag$ inverse $(I_\nfrag - \ext^{(\dom)})^{-1}$ per
domain, which can itself be written symbolically via the matrix
adjugate.
\end{remark}


\subsection{Bilinear Factored Form of Adjacency Frequencies}

The Woodbury identity gives the full path-sum matrix in a form that
makes the adjacency computation practical for any $\ndom, \nfrag$.

\subsubsection{Structure of the path-sum matrix}

Writing $G_\dom = (I - D_\dom)^{-1}$ ($5\nfrag \times 5\nfrag$ per
domain), the Woodbury expansion gives, for states $s_1$ in
$(\dom_1, \frag_1)$ and $s_2$ in $(\dom_2, \frag_2)$:
\[
[(I - \effT)^{-1}]_{s_1 s_2}
= \delta_{\dom_1\dom_2}\,[G_{\dom_1}]_{s_1 s_2}
+ [G_{\dom_1}\,\mathbf{e}_{\dom_1}]_{s_1}^\top\,
  \nonemptytrans_\text{mid}\,K\,
  [\mathbf{s}_{\dom_2}^\top G_{\dom_2}]_{s_2}
\]
where $K = \bigl(I_3 - \sum_\dom \mathbf{s}_\dom^\top G_\dom\, \mathbf{e}_\dom\,
\nonemptytrans_\text{mid}\bigr)^{-1}$
is the $3 \times 3$ Woodbury kernel (push-through form), computed once.

\subsubsection{Bilinear form of structural weights}

The structural weight $C^{\alpha\beta}_{(\dom_1,\frag_1),(\dom_2,\frag_2)}$
decomposes as a within-domain diagonal plus a cross-domain bilinear form:
\[
C^{\alpha\beta}_{(\dom_1,\frag_1),(\dom_2,\frag_2)}
= \delta_{\dom_1\dom_2}\, c^\alpha_{\dom_1\frag_1\frag_2}
+ \mathbf{a}^{\alpha\top}_{\dom_1\frag_1}\, M^{-1}\, \mathbf{b}^{\beta}_{\dom_2\frag_2}
\]
where $\mathbf{a}^\alpha_{\dom\frag}, \mathbf{b}^\beta_{\dom\frag} \in \mathbb{R}^5$
are per-(domain, fragment-type) vectors derived from $G_\dom$,
$\mathbf{e}_\dom$, $\mathbf{s}_\dom$, and the emission-type projections,
and $c^\alpha_{\dom\frag_1\frag_2}$ is the within-domain entry of
$G_\dom$ between fragment-types $\frag_1$ and $\frag_2$ at emission
type $\alpha$ (a $\nfrag \times \nfrag$ matrix per domain rather than a
scalar).

\subsubsection{Closed-form adjacency formula}

Substituting into the adjacency frequency~\eqref{eq:pair-adj} and using
the per-(domain, fragment-type) class-mixture emissions
$\phi^\alpha_{\dom\frag}$ from
Section~\ref{sec:class-mixture-emissions}, the full adjacency table
takes the form:
\begin{equation}\label{eq:bilinear-adj}
f^{\alpha\beta}(\text{chars}_L, \text{chars}_R)
= \underbrace{
  \left[\sum_{\dom,\frag} \mathbf{a}^\alpha_{\dom\frag}\,
        \phi^\alpha_{\dom\frag}(\text{chars}_L)\right]^\top
  M^{-1}
  \left[\sum_{\dom,\frag} \mathbf{b}^\beta_{\dom\frag}\,
        \phi^\beta_{\dom\frag}(\text{chars}_R)\right]
}_{\text{cross-domain: bilinear in two 5-vectors}}
+ \underbrace{
  \sum_\dom \sum_{\frag_1,\frag_2}
      c^\alpha_{\dom\frag_1\frag_2}\,
      \phi^\alpha_{\dom\frag_1}(\text{chars}_L)\,
      \phi^\beta_{\dom\frag_2}(\text{chars}_R)
}_{\text{within-domain: $\nfrag \times \nfrag$ inner sum}}
\end{equation}

This is algebraic closed form: every quantity appearing---$G_\dom$ entries,
$\mathbf{a}^\alpha_{\dom\frag}$, $\mathbf{b}^\beta_{\dom\frag}$,
$c^\alpha_{\dom\frag_1\frag_2}$, $M^{-1}$, and $\phi^\alpha_{\dom\frag}$---is a
rational function of the model parameters (composed with the matrix
exponentials $\exp(\revsub^{(\class)}\evoltime)$).
All inversions reduce to closed-form rational expressions: the only
matrix inversion of dimension exceeding $3 \times 3$ in the entire
distillation pipeline is $(I_\nfrag - \ext^{(\dom)})^{-1}$ per domain
(Section~\ref{sec:within-domain-mixdom}), which is a single
$\nfrag \times \nfrag$ adjugate-over-determinant. The Woodbury
correction $M^{-1}$ is a single $3 \times 3$ inversion regardless of
$\ndom, \nfrag, \nclasses$.

\begin{remark}[Practical computation for $\ndom = \nfrag = \nclasses = 3$]
For $\ndom = 3, \nfrag = 3, \nclasses = 3$:
\begin{enumerate}
\item Precompute $\nclasses = 3$ per-class transition kernels
  $\exp(\revsub^{(\class)}\evoltime)$.
\item Precompute 3 within-domain inverses $G_\dom$ (each
  $15 \times 15$, i.e.\ $5\nfrag = 15$).
\item Precompute 9 pairs $(\mathbf{a}^\alpha_{\dom\frag},
  \mathbf{b}^\beta_{\dom\frag})$ for each $(\alpha,\beta)$ pair, plus
  the within-domain $\nfrag \times \nfrag$ matrices
  $c^\alpha_{\dom\frag_1\frag_2}$.
\item Accumulate $\Sigma = \sum_\dom \mathbf{s}_\dom^\top G_\dom \mathbf{e}_\dom$
  and compute the push-through kernel $K$ (one $3 \times 3$ inversion).
\item For each character entry: evaluate~\eqref{eq:bilinear-adj} by
  summing 9 scaled 3-vectors for the left and right factors of the
  cross-domain term, then a $3 \times 3$ bilinear product, plus a
  within-domain double sum over $(\frag_1, \frag_2)$.
\end{enumerate}
The per-entry cost is $O(\ndom\nfrag + \nclasses)$ multiplications.
\end{remark}


\subsection{Full-Context Distillation: Passthrough Context for Insert and Delete}
\label{sec:full-context}

In a two-tape transducer, the \emph{state} at any point should encode
the most recent character on \emph{each} tape:
\begin{itemize}
\item After Match$(X,Y)$: context is $(X,Y)$ --- both tapes updated.
\item After Insert emitting $Y'$: context is $(X, Y')$ --- ancestor context $X$ unchanged
  (passthrough from prior Match or Delete), descendant updated to $Y'$.
\item After Delete consuming $X'$: context is $(X', Y)$ --- descendant context $Y$ unchanged
  (passthrough from prior Match or Insert), ancestor updated to $X'$.
\end{itemize}

The adjacency frequencies in
Section~\ref{sec:algebraic-distillation} track only partial context
for Insert and Delete states: $\phi^{\ins}_{\dom\frag}(Y)
= \sum_\class \classdist_{\dom\frag\class} \eqm^{(\class)}_Y$
depends only on the descendant character, and
$\phi^{\del}_{\dom\frag}(X) = \sum_\class \classdist_{\dom\frag\class}
\eqm^{(\class)}_X$ depends only on the ancestor character.
This loses information about which (domain, fragment-type, class)
generated the passthrough context.

We now show that the full-context adjacency frequencies---with both
$(X,Y)$ tracked in all states---can be computed in closed algebraic
form, preserving the Woodbury structure.

\subsubsection{Insert chain Green's function}

Define the \emph{Insert chain} as the sub-process restricted to Insert states
$\mathcal{I} = \{\mat\ins_{\dom\frag}, \ins\ins_{\dom\frag}
: \dom \in \ndom, \frag \in \nfrag\}$.
The restricted transition matrix $T^{\mathcal{II}}_{\text{eff}}$
(Insert$\to$Insert transitions within $\effT$) has the same
block-diagonal-plus-low-rank structure as the full $\effT$:
\[
T^{\mathcal{II}}_{\text{eff}} = D^{II} + E^{II}\, \nonemptytrans^{II}_{\text{mid}}\, {S^{II}}^\top
\]
where $D^{II} = \text{diag}(D^{II}_1, \ldots, D^{II}_\ndom)$
with each $D^{II}_\dom$ being the within-domain Insert self-loop block
($2\nfrag \times 2\nfrag$ for $\{\mat\ins_{\dom\frag},
\ins\ins_{\dom\frag}\}$ over fragment-types),
and the cross-domain coupling has rank $\leq 2$ through the Insert rows
of $\nonemptytrans_\bullet$.

The Insert chain Green's function is therefore
\[
G^{II} = (I - T^{\mathcal{II}}_{\text{eff}})^{-1}
\]
computable via Woodbury with a kernel of dimension $\leq 2$.
Each within-domain inverse $(I - D^{II}_\dom)^{-1}$ is a
$2\nfrag \times 2\nfrag$ inversion --- simpler than the
$5\nfrag \times 5\nfrag$ within-domain block of the full system.

An analogous \emph{Delete chain Green's function}
$G^{DD} = (I - T^{\mathcal{DD}}_{\text{eff}})^{-1}$ handles the
descendant passthrough through Delete states, with identical structure.

\subsubsection{Ancestor-conditioned structural weights}

To enter the Insert chain with ancestor context $X$, the process must
have come from a prior Match or Delete state in some $(\dom_0, \frag_0)$
that emitted ancestor $X$.
The match emission satisfies $\sum_Y \phi^{\mat}_{\dom\frag}(X,Y)
= \phi^{\del}_{\dom\frag}(X)$ (row-stochastic per ancestor), so the
entry weight from $(\dom_0, \frag_0)$ with ancestor $X$ is proportional
to $\phi^{\del}_{\dom_0\frag_0}(X)$ for both Match and Delete source
states.

Define the ancestor-conditioned left vector:
\[
\tilde{L}^{\mathcal{I}}_{\dom_1\frag_1}(X) = \sum_{\dom_0,\frag_0}
  \underbrace{(L_{\mat\mat_{\dom_0\frag_0}} + L_{\mat\del_{\dom_0\frag_0}}
              + L_{\del\del_{\dom_0\frag_0}})}_{\lambda_{\dom_0\frag_0}}
  \cdot\; \phi^{\del}_{\dom_0\frag_0}(X)
  \cdot\; \bigl[\effT_{\text{entry}}\, G^{II}\bigr]_{(\dom_0,\frag_0),(\dom_1,\frag_1)}
\]
where $L_s$ are the standard left boundary weights from the full
path-sum computation, $\lambda_{\dom_0\frag_0}$ collects all
ancestor-emitting (Match and Delete) boundary contributions for
$(\dom_0, \frag_0)$, $\effT_{\text{entry}}$ is the entry transition
from ancestor-emitting states $\{\mat\mat, \mat\del, \del\del\}$ into
the Insert chain (i.e.\ the $\{M,D\} \to \{I\}$ block of $\effT$),
and $G^{II}_{(\dom_0,\frag_0),(\dom_1,\frag_1)}$ sums over all
Insert-chain continuations.

Key observation: $\phi^{\del}_{\dom_0\frag_0}(X)$ factors out
multiplicatively from the structural sum, preserving the bilinear
structure.

\subsubsection{Full-context adjacency for Insert-sourced transitions}

The full-context adjacency for Insert$\to$Match is:
\begin{equation}\label{eq:full-im}
f^{\mathcal{I}\mathcal{M}}_{\text{full}}(X, Y, X', Y')
= \sum_{(\dom_1,\frag_1),(\dom_2,\frag_2)}
  \tilde{L}^{\mathcal{I}}_{\dom_1\frag_1}(X)
  \cdot \eta^{\mathcal{IM}}_{(\dom_1,\frag_1),(\dom_2,\frag_2)}
  \cdot \phi^{\ins}_{\dom_1\frag_1}(Y)
  \cdot \phi^{\mat}_{\dom_2\frag_2}(X', Y')
\end{equation}
where $\eta^{\mathcal{IM}}_{(\dom_1,\frag_1),(\dom_2,\frag_2)}
= T_{I_{\dom_1\frag_1}, M_{\dom_2\frag_2}} \cdot R_{M_{\dom_2\frag_2}}$
and $Y$ is the descendant character at the Insert,
$X$ the passthrough ancestor.

The bilinear factored form (cf.~\eqref{eq:bilinear-adj}) generalizes to:
\begin{equation}\label{eq:bilinear-full}
f^{\mathcal{I}\mathcal{M}}_{\text{full}}(X, Y, X', Y')
= \left[\sum_{\dom,\frag} \tilde{\mathbf{a}}^{\mathcal{I}}_{\dom\frag}(X)\,
        \phi^{\ins}_{\dom\frag}(Y) \right]^\top
  K_I
  \left[\sum_{\dom,\frag} \mathbf{b}^{\mathcal{M}}_{\dom\frag}\,
        \phi^{\mat}_{\dom\frag}(X', Y')\right]
  + \sum_{\dom,\frag_1,\frag_2} \text{(diagonal terms)}
\end{equation}
where $\tilde{\mathbf{a}}^{\mathcal{I}}_{\dom\frag}(X)$ are modified
per-(domain, fragment-type) left vectors incorporating the ancestor
context through $\tilde{L}^{\mathcal{I}}_{\dom\frag}(X)$, and $K_I$ is
the Woodbury kernel for the Insert-chain correction.

Similarly, for Insert$\to$Insert with ancestor passthrough:
\[
f^{\mathcal{I}\mathcal{I}}_{\text{full}}(X, Y, X, Y')
= \sum_{(\dom_1,\frag_1),(\dom_2,\frag_2)}
  \tilde{L}^{\mathcal{I}}_{\dom_1\frag_1}(X)
  \cdot \eta^{\mathcal{II}}_{(\dom_1,\frag_1),(\dom_2,\frag_2)}
  \cdot \phi^{\ins}_{\dom_1\frag_1}(Y)
  \cdot \phi^{\ins}_{\dom_2\frag_2}(Y')
\]
where the ancestor $X$ is preserved on both sides (3 effective
character dimensions, not 4).

By symmetry, the Delete-sourced adjacencies with descendant passthrough are:
\begin{equation}\label{eq:full-dm}
f^{\mathcal{D}\mathcal{M}}_{\text{full}}(X, Y, X', Y')
= \sum_{(\dom_1,\frag_1),(\dom_2,\frag_2)}
  \tilde{L}^{\mathcal{D}}_{\dom_1\frag_1}(Y)
  \cdot \eta^{\mathcal{DM}}_{(\dom_1,\frag_1),(\dom_2,\frag_2)}
  \cdot \phi^{\del}_{\dom_1\frag_1}(X)
  \cdot \phi^{\mat}_{\dom_2\frag_2}(X', Y')
\end{equation}
where $\tilde{L}^{\mathcal{D}}_{\dom_1\frag_1}(Y)$ is the
descendant-conditioned left vector for the Delete chain, defined
analogously with $G^{DD}$.

\subsubsection{Computational cost}

The affected adjacency tables gain one character dimension:

\begin{center}
\begin{tabular}{lccc}
\toprule
Tensor & Original & Full-context & Per-entry cost \\
\midrule
$f^{\mathcal{M}\mathcal{M}}$ & $|\alphabet|^4$ & $|\alphabet|^4$ (unchanged) & $O(\ndom\nfrag + \nclasses)$ \\
$f^{\mathcal{M}\mathcal{I}}$ & $|\alphabet|^3$ & $|\alphabet|^3$ (unchanged) & $O(\ndom\nfrag + \nclasses)$ \\
$f^{\mathcal{M}\mathcal{D}}$ & $|\alphabet|^3$ & $|\alphabet|^3$ (unchanged) & $O(\ndom\nfrag + \nclasses)$ \\
$f^{\mathcal{I}\mathcal{M}}$ & $|\alphabet|^3$ & $|\alphabet|^4$ & $O(\ndom\nfrag + \nclasses)$ \\
$f^{\mathcal{I}\mathcal{I}}$ & $|\alphabet|^2$ & $|\alphabet|^3$ ($X$ preserved) & $O(\ndom\nfrag + \nclasses)$ \\
$f^{\mathcal{I}\mathcal{D}}$ & $|\alphabet|^2$ & $|\alphabet|^3$ ($X$ preserved, $Y \to X'$) & $O(\ndom\nfrag + \nclasses)$ \\
$f^{\mathcal{D}\mathcal{M}}$ & $|\alphabet|^3$ & $|\alphabet|^4$ & $O(\ndom\nfrag + \nclasses)$ \\
$f^{\mathcal{D}\mathcal{D}}$ & $|\alphabet|^2$ & $|\alphabet|^3$ ($Y$ preserved) & $O(\ndom\nfrag + \nclasses)$ \\
$f^{\mathcal{D}\mathcal{I}}$ & $|\alphabet|^2$ & $|\alphabet|^3$ ($Y$ preserved, $X \to Y'$) & $O(\ndom\nfrag + \nclasses)$ \\
\bottomrule
\end{tabular}
\end{center}

The per-entry cost is $O(\ndom\nfrag + \nclasses)$ after Woodbury
factoring (cf.\ the same per-entry cost for $f^{\mathcal{M}\mathcal{M}}$
noted in Remark on p.~\pageref{prop:linear}), since the Insert and
Delete chain Green's functions have the same bilinear decomposition.
The total cost for the full adjacency table is
$O((\ndom\nfrag + \nclasses) \cdot |\alphabet|^4)$ --- the same
asymptotic scaling as the existing match-to-match computation.
The Woodbury kernel remains $3 \times 3$ for the full system;
the Insert and Delete chain kernels are $\leq 2 \times 2$.
All quantities remain closed-form rational functions of the model
parameters (composed with the $\nclasses$ rate-matrix exponentials).

\begin{remark}[Why the original formulation lost context]
The emission functions $\phi^{\ins}_{\dom\frag}(Y)$ and
$\phi^{\del}_{\dom\frag}(X)$ are genuinely single-character: the Insert
emission does not depend on the ancestor, and vice versa. The lost
context is not an emission effect but a \emph{transition routing} effect:
the joint posterior $P((\dom,\frag,\class) \mid X, Y)$ depends on both
characters, so the passthrough character informs which
(domain, fragment-type, class) we are in, and hence which transition
probabilities apply. The correction above recovers this information by
conditioning the structural weights on the passthrough character.
\end{remark}


\subsection{Domains versus Fragments versus Classes for Adjacency Capture}
\label{sec:allocation}

The order-1 WFST has $O(|\alphabet|^4)$ free parameters
(from match-to-match transition weights parameterized by context
$(X,Y)$ and next pair $(\anctok',\destok')$), while the MixDom model
has far fewer. We analyse which type of model complexity most
efficiently captures adjacency structure.

\subsubsection{Adjacency tensor rank}

The match-to-match adjacency table $f^{\mathcal{M}\mathcal{M}}(X,Y,X',Y')$
is a sum of $(\ndom\nfrag)^2$ structural-weight terms (in the
(domain, fragment-type)-pair sense), each weighted by a rank-$\nclasses$
emission factor on each side. To approach saturating the WFST capacity
($|\alphabet|^4$ entries), the tensor rank must approach $|\alphabet|^2$;
this is achievable through a combination of $(\ndom, \nfrag, \nclasses)$.

\subsubsection{Domains: independent TKF rates and weights}

Each domain type $\dom$ brings its own $(\alpha_\dom, \beta_\dom,
\gamma_\dom, \kappa_\dom)$, creating genuinely different:
\begin{itemize}
\item $\mat\to\del$ vs.\ $\mat\to\ins$ ratios (different
  $\kappa_\dom(1-\alpha_\dom)$ vs.\ $\beta_\dom$),
\item $\waitm$ vs.\ $\waitd$ transition behaviour (different
  $\beta_\dom$ vs.\ $\gamma_\dom$),
\item overall indel/match balance.
\end{itemize}
These features drive cross-type adjacency diversity in the WFST and
are not reproducible by fragments or classes alone.

\subsubsection{Fragments: vary intra-fragment transition pattern as well as boundary frequency}

The intra-fragment Markov ext matrix
$\ext^{(\dom)}_{\srcfrag\destfrag}$ allows transitions \emph{between
fragment-types within a single fragment}, without invoking the TKF92
new-fragment branch. These intra-fragment $\srcfrag \to \destfrag$
transitions at each emitted site can carry persistent emission-class
context: a chain that prefers one fragment-type tends to stay in that
fragment-type's class mixture, producing richer M/I/D adjacency
patterns than fragments without intra-fragment correlation.
The off-diagonal entries of $\ext^{(\dom)}$ thus contribute genuine new
$\mat/\ins/\del$ adjacency content; setting them to zero
(\texttt{--freeze-offdiag-ext} in the fitter) recovers the
diagonal-extension special case behaviour where fragments only
modulate boundary frequency, and the $\nfrag = 1$ case collapses
fragments to a scalar self-extension per domain.

\subsubsection{Classes: emission-mixture diversity, decoupled from structure}

Each site class $\class$ has its own GTR
$(\exch^{(\class)}, \eqm^{(\class)})$. Classes contribute
\emph{emission} diversity: the rank-$\nclasses$ class-mixture emission
$\phi^{\mat,\ins,\del}_{\dom\frag}$ allows the joint match emission
tensor to deviate from a single GTR's product structure, even within a
single $(d,f)$. Classes do not contribute to the $\mat/\ins/\del$
transition pattern (they do not appear in $\effT$ or in any of the
structural weights $C^{\alpha\beta}$); they affect only the
character-side of each adjacency frequency. Sharing the class pool
across $(d,f)$ via $\classdist$ gives a low-parameter way to enrich
the emission tensor without growing the latent-state space.

\begin{remark}[Parameter allocation]
For adjacency capture in MixDom:
\begin{itemize}[nosep]
\item Domains are essential to enrich the $\mat/\ins/\del$ transition
  pattern via independent TKF rates.
\item Fragments enrich the $\mat/\ins/\del$ structure (via the
  intra-fragment Markov ext) and the boundary frequency
  (via $\notext^{(\dom)}_\srcfrag$); they also support an additional
  layer of emission diversity through $\classdist_{\dom\frag\class}$.
\item Classes enrich emissions only, but cheaply (decoupled from
  $\ndom, \nfrag$): with $\nclasses$ classes, every $(d,f)$ has access
  to the full pool via its Dirichlet $\classdist_{\dom\frag\class}$.
\end{itemize}
The total adjacency-tensor rank scales as $\ndom \nfrag$
(structural-weight pairs) $\times \nclasses^2$ (left/right class mixture);
in practice, allocating the parameter budget across all three axes
gives the most efficient match to the WFST capacity.
\end{remark}


\subsection{Identifiability}
\label{sec:identifiability}

Are the MixDom parameters recoverable from the distilled order-1 models?

\paragraph{Generic identifiability (up to standard ambiguities).}
The distilled WFST provides $O(|\alphabet|^4)$ constraints on the
MixDom parameters. With $\ndom$ domains, $\nfrag$ fragment types per
domain, and $\nclasses$ site classes, there are
$2 + (\ndom-1) + 2\ndom + \ndom\nfrag^2 + \ndom\nfrag(\nclasses-1)
+ \nclasses(|\alphabet|^2 + |\alphabet| - 2)$ free parameters (the
last term accounts for $\nclasses$ symmetric exchangeability matrices
and equilibrium distributions). The system is heavily over-determined
for moderate $(\ndom, \nfrag, \nclasses)$.

The per-domain TKF parameters $(\insrate_\dom\evoltime,
\delrate_\dom\evoltime, \ext^{(\dom)}_{\srcfrag\destfrag})$ are
recoverable from the within-(domain, fragment-type) $\mat/\ins/\del$
transition structure: $\alpha_\dom$ and $\kappa_\dom$ determine
$\delrate_\dom\evoltime$ and $\insrate_\dom/\delrate_\dom$;
$\beta_\dom$ provides a redundant check;
$\ext^{(\dom)}$ separates from $\kappa_\dom$ because intra-fragment
Markov transitions decompose differently from new-fragment transitions
(see Section~\ref{sec:count-restoration} of the Maraschino fitter and
Baum--Welch derivation). The top-level parameters
$(\insrate_\main\evoltime, \delrate_\main\evoltime)$ are identifiable
from the inter-domain transition pattern. The per-class GTRs
$(\exch^{(\class)}, \eqm^{(\class)})$ and per-(domain, fragment-type)
$\classdist_{\dom\frag\class}$ are jointly identifiable up to label
permutation of classes, because the WFST context dependence reveals the
mixture components as the joint posterior shifts across different
$(X, Y)$ contexts.

\paragraph{Unavoidable ambiguities.}
\begin{enumerate}
\item \textbf{Label permutation}: permuting domain labels gives the
  same model ($\ndom!$-fold), permuting fragment-type labels at fixed
  domain gives the same model ($\nfrag!^\ndom$-fold), and permuting
  class labels gives the same model ($\nclasses!$-fold).
\item \textbf{Rate-time confounding}: only products
  $\insrate\evoltime, \delrate\evoltime$ are identifiable from
  pairwise data (a single evolutionary time).
\item \textbf{Class-domain mixing}: when $\nclasses \geq \ndom\nfrag$,
  the class structure can absorb domain-level emission differences,
  creating identifiability issues unless the structural weights
  $C^{\alpha\beta}$ resolve them.
\end{enumerate}

\paragraph{Lossy in distribution, injective in parameters.}
The distilled model captures only pairwise-adjacent correlations;
the full MixDom has higher-order structure (e.g., runs of characters
from the same fragment, intra-fragment Markov correlations between
fragment-types). The distillation map MixDom $\to$ order-1 is
therefore \emph{lossy} for the sequence distribution but generically
\emph{injective} for the parameters: one can recover the MixDom
parameters from the distilled model, even though the distilled model
cannot reproduce all statistics of the MixDom.


\subsection{\texorpdfstring{Scaling to $\ndom, \nfrag, \nclasses$}{Scaling to N, F, C}}
\label{sec:scaling}

\subsubsection{\texorpdfstring{The top-level null closure does not grow with $\ndom, \nfrag, \nclasses$}{The top-level null closure does not grow with N, F, C}}

The null states $\mnull, \inull, \dnull$ are always exactly three,
regardless of the latent-state cardinalities. Their transition submatrix
depends on $\ndom$ only through the scalars
$\emptyseg_0 = \sum_\dom \domdist_\dom(1 - \kappa_\dom)$ and
$\emptyseg_\evoltime = \sum_\dom \domdist_\dom(1 - \kappa_\dom)(1 - \beta_\dom)$;
$\nfrag$ and $\nclasses$ do not appear at this top level.
The null closure and the effective $5 \times 5$ top-level matrix
$\nonemptytrans$ are $O(\ndom)$ to compute but $O(1)$ in matrix
dimension.

\subsubsection{Domain types are drawn i.i.d., not as a Markov chain}

When a new domain begins, its type is drawn independently from
$\domdist$, regardless of the previous domain's type. This i.i.d.\
structure means the cross-domain transitions factor as
$\effT(s, s') = e_\srcdom(s) \times \nonemptytrans_\bullet \times s_\destdom(s')$
(exit $\times$ top-level $\times$ start), and the cross-domain
contribution to the full $5\ndom\nfrag \times 5\ndom\nfrag$ transition
matrix has rank at most 3. Note that this i.i.d.\ property holds at
the \emph{domain} level only; within a fragment, the fragment-type
process is a Markov chain on $\nfrag$ states (see
Section~\ref{sec:nestpairhmm}, Equation~\ref{eq:mixdom_transitions}),
which is encoded in the within-domain block $D_\dom$ and does not
break the cross-domain factorisation.

\subsubsection{\texorpdfstring{Woodbury for general $\ndom, \nfrag$}{Woodbury for general N, F}}

The emitting state space has $5\ndom\nfrag$ states.
The effective transition matrix decomposes as
\[
\effT = D + E\, \nonemptytrans_\bullet\, S^\top
\]
where $D = \text{diag}(D_1, \ldots, D_\ndom)$ with each $D_\dom$ a
$5\nfrag \times 5\nfrag$ within-domain block (now domain-specific and,
for $\nfrag > 1$, with intra-fragment Markov coupling),
$E \in \mathbb{R}^{5\ndom\nfrag \times 3}$ stacks exit vectors projected
to $\{\mat,\ins,\del\}$, and $S \in \mathbb{R}^{5\ndom\nfrag \times 3}$
stacks start vectors projected to $\{\mat,\ins,\del\}$.

The Woodbury identity gives $(I - \effT)^{-1}$ via:
\begin{enumerate}
\item $\ndom$ independent $5\nfrag \times 5\nfrag$ inversions
  $(I - D_\dom)^{-1}$, each evaluated in closed form via the
  Kronecker-plus-rank-3 inner Woodbury of
  Section~\ref{sec:within-domain-mixdom}: one $\nfrag \times \nfrag$
  adjugate $(I_\nfrag - \ext^{(\dom)})^{-1}$ plus a $3 \times 3$ inner
  kernel per domain. For $\nfrag = 1$ this further collapses to a
  $3\times 3$ adjugate (closed-form determinant
  $(1-\ext)(1-\beta)(1-\kappa)$) plus two scalar inversions
  (Section~\ref{sec:within-domain-inversion}). The total cost is
  $O(\ndom \nfrag^3)$ per evaluation.
\item Computing $\sum_\dom \mathbf{s}_\dom^\top G_\dom\, \mathbf{e}_\dom$,
  a $3 \times 3$ matrix built by summing $\ndom$ contributions:
  $O(\ndom \nfrag^2)$ work.
\item One $3 \times 3$ inversion for the Woodbury correction:
  $O(1)$ work.
\end{enumerate}

\begin{proposition}[Linear-in-$\ndom$, cubic-in-$\nfrag$, linear-in-$\nclasses$ scaling]
\label{prop:linear}
The order-1 distillation computation scales as $O(\ndom \nfrag^3)$
in within-domain inversions and $O(\nclasses)$ in per-class GTR
exponentials:
\begin{itemize}
\item $O(\ndom)$ to compute $\emptyseg_0, \emptyseg_\evoltime$ and the
  $5\times 5$ top-level matrix.
\item $O(\ndom \nfrag^3)$ for within-domain path sums.
\item $O(\ndom \nfrag^2)$ to accumulate the $3 \times 3$ Woodbury correction.
\item $O((\ndom\nfrag)^2)$ (domain, fragment-type)-pair terms per
  adjacency entry (or
  $O((\ndom\nfrag)^2 |\alphabet|^d \cdot \nclasses)$ for the full
  adjacency table, with $d \in \{2,3,4\}$ character dimensions
  depending on state type; see Section~\ref{sec:full-context}).
\item $O(\nclasses |\alphabet|^3)$ for per-class transition kernels
  $\exp(\revsub^{(\class)}\evoltime)$.
\end{itemize}
The Woodbury kernel is always $3 \times 3$, regardless of the latent
state cardinalities.
\end{proposition}


\subsection{Summary}

\begin{center}
\begin{tabular}{lcc}
\toprule
Component & Matrix size & Scaling \\
\midrule
Top-level null closure & $3 \times 3$ & $O(1)$ \\
Top-level eff.\ matrix & $5 \times 5$ & $O(\ndom)$ to compute \\
Within-domain inversions & $5\nfrag \times 5\nfrag$ each & $O(\ndom\nfrag^3)$ \\
Woodbury correction & $3 \times 3$ & $O(\ndom\nfrag^2)$ to accumulate \\
Per-class GTR exponentials & $|\alphabet| \times |\alphabet|$ & $O(\nclasses |\alphabet|^3)$ \\
(Domain, fragment-type)-pair adjacency terms & --- & $O((\ndom\nfrag)^2 \nclasses)$ per entry \\
\bottomrule
\end{tabular}
\end{center}

\medskip
The distillation is closed-form for any finite $\ndom, \nfrag, \nclasses$:
all quantities are rational functions of the model parameters
(composed with $\nclasses$ exponentials of rate-times-time products).
The key structural features enabling this are:

\begin{enumerate}
\item \textbf{Top-level null closure is $O(1)$ in dimension.}
  The MixDom null states $\mnull, \inull, \dnull$ are always exactly
  three, regardless of $\ndom, \nfrag, \nclasses$. All dependence on
  $\ndom$ enters through the scalars $\emptyseg_0, \emptyseg_\evoltime$.

\item \textbf{Cross-domain transitions factor.}
  The exit-vector $\times$ top-level $\times$ start-vector structure
  caps the rank of the inter-domain coupling at 3
  (the $\mat, \ins, \del$ top-level states that couple across domains),
  enabling Woodbury reduction to a fixed $3 \times 3$ inversion.

\item \textbf{Domain types are i.i.d., not Markov.}
  If domain type depended on the previous domain (a Markov chain on
  domain types), the cross-domain block would lose its factored
  structure and the Woodbury reduction would fail---the full
  $5\ndom\nfrag \times 5\ndom\nfrag$ inversion would be required,
  scaling as $O((\ndom\nfrag)^3)$. The domain-level i.i.d.\ mixture
  is what keeps the cross-domain part linear in $\ndom$. The
  intra-fragment Markov chain on fragment-types is encoded inside
  $D_\dom$ and does not break this factorisation.

\item \textbf{Class mixture is bilinear in emissions.}
  The class index $\class$ enters every emission as a linear sum
  weighted by $\classdist_{\dom\frag\class}$, never inside the
  structural weights $C^{\alpha\beta}$. This makes the class mixture
  a low-rank emission factor that does not interact with the
  block-diagonal-plus-low-rank structure of the transition matrix.
\end{enumerate}

\ifdefined\INCLUDEDFROMTKF\else
\end{document}
\fi
